\scanpage{92}
% archive foliation: 91
\begin{center}- 137 -\end{center}
\hfill n\textsuperscript{o} 316
\begin{center}- \uline{Exercices sur le §3} -\end{center}

1) - Soient $g$ une algèbre de Lie, $\rho$ une représentation de $g$ dans un $K$-module $M$, $\sigma$ la représentation associée de $g$ dans l'algèbre tensorielle de $M$. Montrer que le sous-module des tenseurs symétriques et le sous-module des tenseurs antisymétriques sont stables pour $\sigma$.

2) - Soient $g$ une algèbre de Lie, $\rho$ une représentation de $g$ dans un $K$-module $M$, $\sigma$ la représentation associée de $g$ dans $Q=\mathcal L(M,M;M)$. Pour qu'un $f\in Q$ soit invariant pour $\sigma$, il faut et il suffit que les $x_M$ soient des dérivations de $M$ muni de la multiplication définie par $f$.

3) - Soit $V$ un espace vectoriel de dimension finie sur un corps parfait $K$.

a) La représentation identique de $\mathrm{gl}(V)$ définit canoniquement des représentations de $\mathrm{gl}(V)$ dans $\overset{p}{\bigotimes}V$ et $\overset{q}{\bigotimes}V^*$, donc une représentation $u\longrightarrow u_q^p$ de $\mathrm{gl}(V)$ dans $V_q^p=(\overset{p}{\bigotimes}V)\otimes(\overset{q}{\bigotimes}V^*)$. Montrer que $u_q^p$ est semi-simple si $u$ l'est. (Se ramener, par extension du corps de base, au cas où $K$ est algébriquement clos). Montrer que $u_q^p$ est nilpotent si $u$ l'est. En déduire que, si $s$ et $n$ sont les composantes semi-simple et nilpotente de $u$, $s_q^p$ et $n_q^p$ sont les composantes semi-simple et nilpotente de $u_q^p$.

b) Déduire de a) que, si un élément de $V_q^p$ est annulé par $u_q^p$, il est annulé par $s_q^p$ et $n_q^p$.

c) Déduire de b) et de l'exerc.2 que, si $V$ est muni d'une structure d'algèbre non nécessairement associative, et si $u$ est une
